\documentclass[10pt]{article}
\usepackage[margin=20mm]{geometry}
\usepackage{amsmath,amssymb,mathtools,fontspec,unicode-math,xurl,hyperref}
\IfFontExistsTF{DejaVu Serif}{\setmainfont{DejaVu Serif}}{\setmainfont{Latin Modern Roman}}
\setmathfont{latinmodern-math.otf}
\setlength{\parindent}{0pt}\setlength{\parskip}{6pt}\sloppy
\title{Technical Review of the Compactness Hypotheses and Final Tangent-Flow Iteration}
\author{AI-assisted technical review}
\date{}
\begin{document}
\maketitle
Abstract: Lemma 3.1, as a statement about arbitrary sequences of curves, permits escape to infinity and therefore does not ensure the nonempty shrinker asserted in its conclusion. Translated unit circles give an explicit example. The actual ancient-flow application has additional entropy and monotonicity bounds that exclude this example; the issue is classified Minor with that impact boundary. Independent checks of the Gaussian modes, a rotation norm estimate, and the final geometric iteration are consistent with the claimed convergence, conditional on the upstream geometric and analytic estimates. This initial review does not certify the full uniqueness theorem. Scientific adjudication and independent acceptance remain pending.\par
\section*{Source and scope}
M25-0377 is Kyeongsu Choi, Dong-Hwi Seo, Wei-Bo Su and Kai-Wei Zhao, Uniqueness of Tangent Flows at Infinity for Finite-Entropy Shortening Curves, Geometric and Functional Analysis 35 (2025), 408–462, DOI 10.1007/s00039-025-00705-1 [source]. The assigned source is the 55-page version-of-record PDF. Printed page \(p\) is physical page \(p-407\). The current publisher HTML was consulted for identity and the exact scope of Lemma 3.1; no separate preprint comparison is asserted.\par
The main setting is a smooth ancient embedded planar curve-shortening flow, with rescaling \(\overline M_\tau=(-t)^{-1/2}M_t\), \(\tau=-\log(-t)\), and finite entropy, the supremum of Gaussian length over all centers and positive scales and all times. Section 2.3 explicitly works with complete curves; the weak flows are integral, unit-regular and cyclic. Incomplete curve fragments are therefore not used as counterexamples.\par
The review checks the nonvanishing boundary of Lemma 3.1, its application to rough convergence, and the algebra and final iteration in Section 8. The detailed low-entropy classification, finger and vertex geometry, full proof of the graphical-radius estimate, and all external regularity theorems remain outside independent certification.\par
\section*{Problem 1: the general compactness lemma allows an empty limit}
Classification: Minor. Location: Lemma 3.1, printed p. 420 (physical page 13), especially the passage from local compactness to a nonempty line or circle. Its application is Theorem 3.3, pp. 421–423.\par
Equation reference: defect. The lemma assumes only that \(\overline M_j\) are connected smooth embedded curves with \(F(\overline M_j)<\Lambda\), and \[
 \int_{\overline M_j}
 \left|\boldsymbol\kappa+\frac{x^\perp}{2}\right|^2
 e^{-|x|^2/4}\,d\mathcal H^1\longrightarrow0.\] It concludes subsequential local \(C^{1,\alpha}\) convergence to a line through the origin or a radius-\(\sqrt2\) circle, possibly with multiplicity.\par
For integers \(j\ge3\), take the complete embedded unit circle \[\overline M_j=\{(j,0)+(\cos s,\sin s):0\le s<2\pi\}.\] These curves are connected and smooth, with \(|\boldsymbol\kappa|=1\), total length \(2\pi\), and \(j-1\le |x|\le j+1\). Consequently, \[\begin{aligned}
 F(\overline M_j)&\le\sqrt\pi\,e^{-(j-1)^2/4}<1,\\
 \int_{\overline M_j}\left|\boldsymbol\kappa+\frac{x^\perp}{2}\right|^2
 e^{-|x|^2/4}\,d\mathcal H^1
 &\le2\pi\left(1+\frac{j+1}{2}\right)^2e^{-(j-1)^2/4}
 \longrightarrow0.\end{aligned}\] Thus the lemma’s hypotheses hold with \(\Lambda=1\). For every fixed ball, however, the curves are eventually absent. Every local measure limit is zero; there is no nonempty line or circle to which a subsequence converges. The construction even has uniformly bounded full entropy, by translation invariance. A lower bound preventing loss of all mass is essential.\par
\subsection*{Why the ancient-flow application is different}
Equation reference: anchor, tail. This is not a counterexample to Theorem 1.1. In that application, fix a nonempty time slice at \(\tau_*\). Huisken monotonicity gives \[
 F(\overline M_\tau)\ge F(\overline M_{\tau_*})=:c>0
 \qquad(\tau\le\tau_*).\] The full entropy bound \(\operatorname{Ent}(\overline M_\tau)\le\Lambda\) also implies uniform linear length growth. Taking Gaussian scale \(R^2\) yields \[\mathcal H^1(\overline M_\tau\cap B_R)
 \le \sqrt{4\pi}\,e^{1/4}\Lambda R .\] Therefore, for integer \(N\ge1\), the Gaussian mass outside \(B_N\) is at most \[
 \frac1{\sqrt{4\pi}}\int_{\overline M_\tau\setminus B_N}e^{-|x|^2/4}\,d\mathcal H^1
 \le e^{1/4}\Lambda\sum_{k=N}^{\infty}(k+1)e^{-k^2/4}
 \longrightarrow0\] uniformly in \(\tau\). Choose \(N\) making this tail smaller than \(c/2\). Equations [anchor] (\#anchor)–[tail] (\#tail) then retain positive mass in a fixed ball.\par
The local repair is to add an explicit nonvanishing condition to the general lemma, or allow an empty limit and exclude it in the application using the bounds above. The necessary nonvanishing is available under the paper’s ancient-flow assumptions. This verifies the repair of the exhibited escape mechanism, not every remaining compactness or shrinker-classification step. Because the main application supplies what is missing, no substantive failure of the main uniqueness conclusion is established; the issue is Minor.\par
\section*{Checks of the final analytic mechanism}
\subsection*{Gaussian modes and stable energy}
On printed p. 450 the graph equation is \[u_\tau=\mathcal Lu-\frac{u_y^2u_{yy}}{1+u_y^2},
 \qquad \mathcal L=\partial_y^2-\frac y2\partial_y+\frac12,
 \qquad d\gamma=\frac{e^{-y^2/4}}{\sqrt{4\pi}}\,dy.\] Direct calculation gives the normalized modes \[\varphi_1=1,\quad \varphi_2=y/\sqrt2,\quad
 \varphi_3=(y^2-2)/\sqrt8,\qquad
 \mathcal L\varphi_1=\tfrac12\varphi_1,\quad
 \mathcal L\varphi_2=0,\quad
 \mathcal L\varphi_3=-\tfrac12\varphi_3.\] The Gaussian moments are \(E[y^2]=2\), \(E[y^4]=12\), so these normalizations agree with the source. The positive number attached to the stable mode is its decay rate, not a positive eigenvalue of \(\mathcal L\).\par
For a stable projection \(w=P_-\widehat u\), \[\langle w,\mathcal Lw\rangle
 =-\|w_y\|_{\mathcal H}^2+\tfrac12\|w\|_{\mathcal H}^2,
 \qquad
 \langle w,\mathcal Lw\rangle\le-\tfrac12\|w\|_{\mathcal H}^2.\] Splitting twice this inner product into \(3/2\) and \(1/2\) multiples and using Young’s inequality for \(2\langle P_-E,w\rangle\) gives \[\frac d{d\tau}\|w\|_{\mathcal H}^2
 \le-\tfrac12\|w_y\|_{\mathcal H}^2
    -\tfrac14\|w\|_{\mathcal H}^2+4\|P_-E\|_{\mathcal H}^2.\] This verifies the coercive arithmetic underlying (68), conditional on the error estimates in Lemma 8.4. The entire nonlinear projection estimate is not certified here.\par
\subsection*{Rotation preserves the needed \(L^2\) control}
Proposition 8.8, printed p. 456, passes from a nearly linear graph to a rotated graph. A direct change of variable supplies the norm control needed for its regularity step. For rotation by \(-\theta\), write \[x=y\cos\theta+u(y)\sin\theta,\qquad
 v(x)=u(y)\cos\theta-y\sin\theta
      =\cos\theta\,[u(y)-(\tan\theta)y].\] When \(|\theta|\) and \(|u_y|\) are small, \(1/2\le dx/dy\le2\); the interior interval \([-6,6]\) maps back into \([-8,8]\). Hence \[\|v\|_{L^2([-6,6])}\le\sqrt2\,
 \|u-(\tan\theta)y\|_{L^2([-8,8])}.\] This closes the norm conversion directly, without inferring an \(L^2\) estimate from an \(L^1\) estimate alone. The cited parabolic local boundedness and Schauder estimates are retained as external inputs.\par
\subsection*{Geometric iteration and exponential rate}
Assume Lemma 8.9 with fixed \(L,C_3,\delta_0\), and coherent choices of the improved rotations. At \(\tau_k=\tau_1-(k-1)L\), the flatness is \(2^{1-k}\delta_0\), and the successive rotations satisfy \[|S_{k-1}-S_k|\le 2^{2-k}C_3\delta_0\quad(k\ge2).\] The exact geometric sum gives \[|S_\infty-S_k|
 \le C_3\delta_0\sum_{j=k+1}^{\infty}2^{2-j}
 =2^{2-k}C_3\delta_0.\] Thus the limit rotation exists. The time windows of length \(4L\), spaced by \(L\), cover the backward time ray. Since \(2^{-k}=C e^{(\log2/L)\tau_k}\), the estimates give exponential decay in fixed coordinates.\par
Choose \(0<5\delta<\log2/L\). The bound \(Ce^{5\delta\tau}\) yields a graphical radius of order \(e^{-5\delta\tau/4}\) through \(\rho^{-4}=\max_i\|u_i\|_{C^2([-1,1])}\) in Theorem 7.3; this eventually exceeds \(e^{-\delta\tau}\). For each fixed \(R\), the remaining prefactor is absorbed by \(Ce^{5\delta\tau}\le e^{4\delta\tau}\) for sufficiently negative \(\tau\), after the stated interior regularity step. These exponent conversions agree with Theorem 1.2. This is a conditional check of the final implication, not a new proof of the upstream flatness improvement.\par
\section*{Coverage and verification}
An independently written Python calculation checked exact polynomial eigenvalues, Gaussian normalization, the stable-energy coefficients and geometric sums. It also evaluated the explicit escape bounds in logarithmic form. The script exited successfully; the limiting arguments are established analytically above. No author code was executed.\par
Item | Status | Decisive check | Boundary\par
Lemma 3.1 | Minor | Escaping unit circles | Application has mass anchor\par
Spectral arithmetic | Consistent | Exact modes/energy | Nonlinear estimates assumed\par
Final iteration | Consistent | Rotation tail and exponents | Lemma 8.9 assumed\par
Full theorem | Unverified | Targeted review only | Geometry/classification open to review\par
Sections 4–7’s complete geometric arguments, exact critical-point counts, weak-flow classification, and external reference proofs were not independently verified. Appendix A’s entropy obstruction to a nonproper complete curve was read and is consistent with the distinction between completeness and an arbitrary incomplete fragment.\par
Bounded correction/version searches did not identify a correction; this is not evidence of exhaustive absence. Publisher HTML retains the same Lemma 3.1 hypotheses. Acquisition and optional-resource search limitations remain in the companion record. The report and evidence-page inspection records bind actual PNG displays to exact source/PDF hashes; compiler, dependency and inventory records bind the final report to its TeX. This is an initial review with scientific decision pending and no human reviews.\par
9 K. Choi, D.-H. Seo, W.-B. Su and K.-W. Zhao, Uniqueness of Tangent Flows at Infinity for Finite-Entropy Shortening Curves, Geom. Funct. Anal. 35 (2025), 408–462. Assigned version-of-record PDF. Springer Nature, article HTML, https://link.springer.com/article/10.1007/s00039-025-00705-1, consulted for identity and Lemma 3.1 wording.\par
\end{document}
